\subsection{Formalización Matemática del Espacio Latente Multimodal y Representaciones Matryoshka}
\label{subsec:formulacion_matematica_mrl}

La arquitectura propuesta proyecta datos no estructurados de naturalezas heterogéneas (vídeo, audio y texto) en un espacio métrico continuo unificado $\mathcal{S} \subseteq \mathbb{R}^{D_b}$, donde $D_b$ representa la dimensionalidad nativa del motor de vectorización $b$ seleccionado bajo el protocolo común \texttt{MultimodalVectorizer} ($D_{\text{gemini}} = 3072$ para \textit{Gemini Embedding 2} \cite{gemini_emb2_docs}, $D_{\text{qwen}} = 2048$ para \textit{Qwen3-VL-Embedding-2B} \cite{qwen3_vl_emb2026, hf_qwen3_2b} y $D_{\text{gemma}} = 768$ para \textit{EmbeddingGemma 2} \cite{hf_embeddinggemma2}).

\subsubsection{Espacio de Ingesta y Proyección Omnimodal}
Sea un flujo de vídeo segmentado temporalmente en una secuencia ordenada de fragmentos visuales $\{C_1, C_2, \dots, C_N\}$. Cada fragmento $C_i$ queda definido por una tupla multimodal:
\begin{equation}
    C_i = \left( \mathcal{F}_i, \tau_i \right)
\end{equation}
donde $\mathcal{F}_i = \{f_{i,1}, f_{i,2}, f_{i,3}\} \subset \mathbb{R}^{H \times W \times 3}$ corresponde al conjunto ordenado de tres fotogramas clave representativos del plano visual (fotogramas inicial, intermedio y final), y $\tau_i \in \mathcal{T}$ denota la transcripción fonética continua alineada temporalmente con dicho intervalo.

La función de incrustación omnimodal $\mathcal{E}: \left( (\mathbb{R}^{H \times W \times 3})^3 \times \mathcal{T} \right) \to \mathbb{R}^D$ transforma la tupla $C_i$ en un vector denso latente $\mathbf{v}_i$:
\begin{equation}
    \mathbf{v}_i = \mathcal{E}_b(\mathcal{F}_i, \tau_i), \quad \|\mathbf{v}_i\|_2 = 1, \quad \mathbf{v}_i \in \mathbb{S}^{D_b-1}
\end{equation}
donde $\mathbb{S}^{D_b-1}$ denota la hiperesfera unitaria en $\mathbb{R}^{D_b}$. De manera análoga, una consulta multimodal arbitraria $\mathcal{Q}$ (compuesta por cualquier combinación no vacía de texto libre $t_q$, imágenes de referencia $\{I_{q,j}\}$ y señales auditivas transcritas $\tau_q$) es proyectada por la misma función en el vector unitario $\mathbf{q} \in \mathbb{S}^{D_b-1}$:
\begin{equation}
    \mathbf{q} = \mathcal{E}_b(\mathcal{Q}), \quad \|\mathbf{q}\|_2 = 1
\end{equation}

\subsubsection{Representaciones Anidadas Matryoshka (MRL)}
A fin de optimizar el consumo de memoria en memoria principal y acelerar las operaciones de búsqueda aproximada sobre grafos jerárquicos (HNSW), el sistema adopta el paradigma \textit{Matryoshka Representation Learning} (MRL) \cite{kusupati2022matryoshka}. MRL entrena el codificador imponiendo una función de pérdida multi-escala que fuerza a los primeros $d$ componentes del vector a concentrar la máxima varianza y capacidad discriminativa semántica.

Formalmente, para cualquier dimensión objetivo admisible $d \in \mathcal{D}_{\text{MRL}}(b)$ ($\mathcal{D}_{\text{MRL}}(\text{gemini}) = \{3072, 1536, 768\}$, $\mathcal{D}_{\text{MRL}}(\text{qwen}) = \{2048, 1536, 768\}$ y $\mathcal{D}_{\text{MRL}}(\text{gemma}) = \{768, 512, 256\}$), se define el operador de truncamiento $\pi_d: \mathbb{R}^{D_b} \to \mathbb{R}^d$ que selecciona el prefijo inicial del vector:
\begin{equation}
    \pi_d(\mathbf{v}) = (v_1, v_2, \dots, v_d)^T
\end{equation}

Dado que la operación de corte altera la norma euclídea del vector ($\|\pi_d(\mathbf{v})\|_2 \le 1$), es imperativo aplicar una proyección normalizadora $L_2$ para preservar la coherencia geométrica de la similitud del coseno:
\begin{equation}
    \mathcal{P}_d(\mathbf{v}) = \begin{cases}
        \dfrac{\pi_d(\mathbf{v})}{\|\pi_d(\mathbf{v})\|_2}, & \text{si } \|\pi_d(\mathbf{v})\|_2 > 0 \\
        \pi_d(\mathbf{v}), & \text{en caso contrario}
    \end{cases}
\end{equation}

Bajo esta transformación, la métrica de similitud entre una consulta $\mathbf{q}$ y un candidato $\mathbf{v}$ en el subespacio $d$-dimensional equivale al producto escalar estándar:
\begin{equation}
    S_d(\mathbf{q}, \mathbf{v}) = \langle \mathcal{P}_d(\mathbf{q}), \mathcal{P}_d(\mathbf{v}) \rangle = \sum_{j=1}^d \mathcal{P}_d(\mathbf{q})_j \cdot \mathcal{P}_d(\mathbf{v})_j
\end{equation}

Esta formulación garantiza que un único índice vectorial pueda ser consultado a múltiples resoluciones de granularidad dimensional sin necesidad de re-entrenar ni re-vectorizar el corpus audiovisual.
